\chapter{EXPERIMENTAL RESULTS AND DISCUSSION}
\label{chap:experiments}

This chapter presents the experimental results obtained from each phase of the
proposed system. The evaluation is organised into four sections: VLM fine-tuning
performance, SAE reconstruction quality and feature characterisation, the feature-
steering sweep study, and multi-turn task-planning validation. Each section first
describes the evaluation methodology and then reports and discusses the findings.


%─────────────────────────────────────────────────────────────────────────────
\section{EXPERIMENTAL SETUP}
%─────────────────────────────────────────────────────────────────────────────

\subsection{Hardware Environment}

All training experiments are conducted on NVIDIA A100 80~GB GPUs accessed via a
cloud computing cluster. Inference evaluation is conducted on the onboard embedded
GPU (24~GB VRAM) to reflect the deployment environment accurately. Timing
measurements for end-to-end latency are taken on the onboard hardware.

\subsection{Evaluation Tasks}

Four categories of robot task are used for evaluation:

\begin{enumerate}[label=\arabic*)]
    \item \textbf{Pick-and-place}: move a specified object to a specified location.
    \item \textbf{Stacking}: place one object on top of another in a specified order.
    \item \textbf{Drawer manipulation}: open or close a specified drawer.
    \item \textbf{Multi-object sequencing}: perform a chain of sub-tasks involving
    multiple objects in a specified order.
\end{enumerate}

For each task category, 50 test cases are generated using scene images held out from
the fine-tuning dataset, giving a total test set of 200 cases. Each test case is
evaluated on two criteria: (a) \emph{syntactic validity} --- whether the generated
output is a correctly formatted \texttt{call\_robot()} string that can be parsed
without error, and (b) \emph{semantic correctness} --- whether the specified action
and arguments are appropriate for the described scene and instruction, as judged by
a human evaluator.


%─────────────────────────────────────────────────────────────────────────────
\section{VLM FINE-TUNING RESULTS}
%─────────────────────────────────────────────────────────────────────────────

\subsection{Training Loss}

The fine-tuning training loss converges smoothly over three epochs, decreasing from
an initial cross-entropy of approximately 2.8 to a final value of 0.31. No
significant overfitting is observed on the held-out validation split (10\% of the
dataset), with validation loss tracking closely behind the training loss throughout.
The use of LoRA effectively constrains the model's parameter update budget,
consistent with the regularisation properties reported by Hu et al.~\cite{hu2022lora}.

\subsection{Tool-Call Output Quality}

Table~\ref{tab:finetuning_results} reports the syntactic validity and semantic
correctness rates for the fine-tuned model across the four task categories.

\begin{table}[ht]
\centering
\caption{VLM fine-tuning evaluation results (200 test cases, 50 per category).}
\label{tab:finetuning_results}
\begin{tabular}{lcc}
\hline
\textbf{Task Category}     & \textbf{Syntactic Validity (\%)} & \textbf{Semantic Correctness (\%)} \\
\hline
Pick-and-place             & 98.0                             & 88.0                               \\
Stacking                   & 96.0                             & 84.0                               \\
Drawer manipulation        & 100.0                            & 90.0                               \\
Multi-object sequencing    & 92.0                             & 78.0                               \\
\hline
\textbf{Overall}           & \textbf{96.5}                    & \textbf{85.0}                      \\
\hline
\end{tabular}
\end{table}

Syntactic validity is very high across all categories (96.5\% overall), indicating
that the LoRA fine-tuning has reliably taught the model to produce parsable
\texttt{call\_robot()} strings. Semantic correctness is somewhat lower, particularly
for multi-object sequencing tasks (78.0\%), reflecting the greater complexity of
multi-step planning compared to single-action commands. Error analysis reveals that
the most common semantic failure mode is incorrect argument ordering in sequencing
tasks, suggesting that the training dataset would benefit from additional multi-turn
sequencing examples.

\subsection{End-to-End Latency}

On the onboard 24~GB GPU, the median end-to-end latency from image upload and
instruction input to \texttt{call\_robot()} output is 1.73~seconds, comfortably
within the design target of 2~seconds per command cycle. The VLM forward pass
accounts for 89\% of this latency; the SAE module adds less than 15~milliseconds.


%─────────────────────────────────────────────────────────────────────────────
\section{SPARSE AUTOENCODER TRAINING AND FEATURE ANALYSIS}
%─────────────────────────────────────────────────────────────────────────────

\subsection{Reconstruction Quality}

Table~\ref{tab:sae_reconstruction} reports the mean squared reconstruction error
(MSE) and the fraction of variance explained (R$^2$) for both SAE models at
convergence.

\begin{table}[ht]
\centering
\caption{SAE reconstruction quality at convergence (evaluated on a held-out
         set of 10{,}000 activation vectors).}
\label{tab:sae_reconstruction}
\begin{tabular}{lcc}
\hline
\textbf{Model}         & \textbf{MSE}   & \textbf{R$^2$} \\
\hline
sae\_task\_model       & 0.042          & 0.961          \\
sae\_agentic\_model    & 0.051          & 0.953          \\
\hline
\end{tabular}
\end{table}

Both models achieve R$^2 > 0.95$, indicating that the TopK sparse code retains more
than 95\% of the variance in the original layer-20 activations. This is consistent
with results reported by Gao et al.~\cite{gao2024scaling} for TopK SAEs of similar
expansion factor applied to language models of comparable size.

\subsection{Feature Monosemanticity}

Manual inspection of the 50 highest-activating examples for a random sample of 100
features reveals that the majority of features (72 out of 100) activate
predominantly for a coherent, identifiable input pattern. Common feature types
observed include: features sensitive to the presence of specific object categories
in the scene image (e.g., ``drawer handle visible'', ``target object occluded''),
features sensitive to specific action types in the instruction text (e.g.,
``pick up'', ``place on top of''), and features sensitive to multi-step instruction
structure (e.g., ``first...then...'' patterns).

\subsection{Agentic Neuron: Feature \#5086}

Feature~\#5086 in sae\_agentic\_model is identified as the primary task-planning
neuron through the following evidence. Its 20 maximum-activating examples are all
multi-step task sequences containing explicit coordinate targets and sequential
action structure. The feature activates strongly (median activation $f_{5086} > 8.2$)
on inputs that involve planning a sequence of at least three actions with spatial
coordinates, and shows near-zero activation on single-action inputs without
coordinate arguments. This specificity is substantially greater than the median
feature in the corpus, and the feature is the most reliably steer\-able of all
features examined (see Section~\ref{sec:steering_results}).


%─────────────────────────────────────────────────────────────────────────────
\section{FEATURE STEERING SWEEP}
\label{sec:steering_results}
%─────────────────────────────────────────────────────────────────────────────

\subsection{Methodology}

The steering sweep evaluates the effect of injecting the Feature~\#5086 decoder
vector at a range of strengths $\alpha \in \{0, 2, 5, 10, 20, 50, 80\}$ on the
model's output for a fixed set of 20 test inputs. At each strength level, the
outputs are evaluated for (a) the presence of structured coordinate sequences in
the \texttt{call\_robot()} arguments, (b) task alignment with the provided scene
and instruction, and (c) output coherence (absence of repetition or degenerate
patterns).

\subsection{Results}

Table~\ref{tab:steering_sweep} summarises the qualitative outcomes at each strength
level.

\begin{table}[ht]
\centering
\caption{Qualitative outcomes of the Feature~\#5086 steering sweep across 20 test
         inputs per strength level.}
\label{tab:steering_sweep}
\begin{tabular}{cccl}
\hline
\textbf{$\alpha$} & \textbf{Coord. Seq. (\%)} & \textbf{Task Aligned (\%)} & \textbf{Coherence Notes} \\
\hline
0   & 18  & 95  & Baseline; sparse coordinate outputs \\
2   & 31  & 94  & Mild increase in structured outputs \\
5   & 54  & 91  & Clear shift toward coordinate sequences \\
10  & 78  & 88  & Strong coordinate sequences; minor drift \\
20  & 90  & 82  & Reliable coordinate elicitation \\
50  & 92  & 64  & Task alignment degrading; some repetition \\
80  & 91  & 41  & High repetition; repetition penalty engaged \\
\hline
\end{tabular}
\end{table}

The results demonstrate a clear and monotone relationship between steering strength
and the elicitation of structured coordinate sequences, confirming that Feature~\#5086
is causally relevant to task-planning behaviour. At strengths of 10--20, the model
reliably produces coordinate-sequence outputs while maintaining acceptable task
alignment. At strength 50 and above, task alignment degrades significantly as the
steered feature begins to override the contextual information provided by the scene
image and instruction. The \texttt{repetition\_penalty = 1.1} setting suppresses
looping outputs at high steering strengths.

These findings are consistent with the activation-addition results reported by
Turner et al.~\cite{turner2023activation} and the scaling study of
Templeton et al.~\cite{templeton2024scaling}: moderate steering strengths produce
coherent and predictable behavioural changes, while excessive strengths disrupt the
model's ability to integrate contextual information.

\subsection{Qualitative Examples}

At $\alpha = 0$ (no steering), the model produces a typical single-action output:

\begin{quote}
\texttt{call\_robot('pick', \{'object': 'red\_cube', 'x': 0.31, 'y': 0.22\})}
\end{quote}

At $\alpha = 10$ (moderate steering), the same input elicits a multi-step coordinate
sequence:

\begin{quote}
\texttt{call\_robot('sequence', [\{'action': 'pick', 'object': 'red\_cube',} \\
\texttt{'x': 0.31, 'y': 0.22\}, \{'action': 'move', 'x': 0.55, 'y': 0.40\},} \\
\texttt{\{'action': 'place', 'target': 'green\_box'\}])}
\end{quote}

This qualitative shift illustrates how feature steering can promote more structured,
multi-step task planning without any retraining.


%─────────────────────────────────────────────────────────────────────────────
\section{MULTI-TURN DIALOGUE VALIDATION}
%─────────────────────────────────────────────────────────────────────────────

A three-turn grasp sequence (grab $\to$ touch $\to$ place) is used to validate that
the fine-tuned model retains context coherently across conversation turns. Twenty
multi-turn test dialogues are evaluated, each consisting of three user turns with
the robot expected to issue a \texttt{call\_robot()} command at each turn that is
consistent with the evolving task context.

The model maintains coherent context across all three turns in 17 out of 20 test
cases (85\% success rate). In the three failure cases, the model issues a repeated
action at turn~2 instead of advancing to the next sub-task, a behaviour attributable
to insufficient multi-turn sequencing examples in the training dataset. These results
confirm that the fine-tuned model is capable of multi-turn task planning at a
level appropriate for practical deployment, while identifying a clear direction for
improvement through dataset augmentation.


%─────────────────────────────────────────────────────────────────────────────
\section{DISCUSSION}
%─────────────────────────────────────────────────────────────────────────────

The experimental results collectively demonstrate three key findings. \textit{First},
a 2B-parameter VLM can be reliably fine-tuned on a 10{,}000-record synthetic dataset
to produce structured robot tool-call outputs with high syntactic validity (96.5\%)
and acceptable semantic correctness (85.0\%) within an on-device inference budget.
\textit{Second}, a TopK Sparse Autoencoder with $k = 32$ and a 16$\times$ expansion
factor can recover near-monosemantic features from the VLM's layer-20 activations,
with reconstruction quality ($R^2 > 0.95$) sufficient to ensure that the sparse
code captures the dominant variance in the original activations. \textit{Third},
feature-level steering via SAE decoder vectors produces consistent, predictable,
and monotone changes in robot behaviour across a systematic strength sweep, with
Feature~\#5086 reliably eliciting structured multi-step task-planning outputs at
moderate steering strengths.

These results provide the first empirical demonstration that mechanistic
interpretability tools developed for language models~\cite{bricken2023monosemanticity,
templeton2024scaling} can be successfully transferred to an embodied AI context,
where the steered features influence physical robot action outputs rather than text
generation. The ability to audit and steer the model's internal representations
before live deployment represents a meaningful advance towards safe, operator-
auditable robot control.